\section{The model and elementary reductions}\label{sec:preliminaries}
We first specify the full strategy space. The later reductions will select
equivalent candidates from this space, without assuming that the five
players share a rule.
Let $Q_q=\{1,\ldots,q\}$. The vertices of $C_5$ are $1,\ldots,5$ in cyclic
order; vertex indices are read modulo~5. Set
\[
\calC_q=\{c\in Q_q^5:c_i\ne c_{i+1}\text{ for every }i\}.
\]
In particular, $c_5\ne c_1$ is required. Nonadjacent vertices may have the
same color, and a member of $\calC_q$ need not use all $q$ colors. A
strategy is a tuple $f=(f_1,\ldots,f_5)$, where
\[
f_i:Q_q^2\longrightarrow Q_q.
\]
The first and second arguments are the left and right neighbor colors,
respectively. Winning means
\begin{equation}\label{eq:winning}
\forall c\in\calC_q\quad
\exists i\in\{1,\ldots,5\}\quad
f_i(c_{i-1},c_{i+1})=c_i.
\end{equation}
There is no communication after the colors have been assigned.
We use
\[
\HG(C_5)=\sup\{q\in\mathbb Z_{\geq2}:\text{a winning strategy with $q$ colors exists}\}.
\]
The lower and upper bounds below make this supremum a finite maximum.

\begin{lemma}[Restriction of the palette]\label{lem:palette}
If a winning strategy exists with $q$ colors and $2\leq r\leq q$, then a
winning strategy exists with $r$ colors.
\end{lemma}
\begin{proof}
Choose an $r$-element subset $R$ of the palette and restrict attention to
configurations with colors in $R$. On observations in $R^2$, replace any
answer outside $R$ by an arbitrary element of $R$. Such an original answer
could not have been correct on a configuration under consideration, so no
previously correct guess is lost. Relabel $R$ as $Q_r$.
\end{proof}

\begin{proposition}[The four-color lower bound]\label{prop:four}
There is a winning strategy on $C_5$ with four colors.
\end{proposition}
\begin{proof}
Use the path strategy of~\cite[Lemma~4.1]{proper}. Identify the colors with
$\F_2^2$ and write the colors at three consecutive vertices as $a,b,c$.
Their guesses are, respectively,
\[
(b_1+1,b_2),\qquad(c_1+1,a_2+1),\qquad(b_1,b_2+1).
\]
All coordinates and additions in this construction are in $\F_2$.
The end vertices ignore their other neighbor. Suppose both end guesses
fail. Properness excludes $a=(b_1,b_2)$, and failure of the first guess
excludes $a=(b_1+1,b_2)$. These are exactly the two colors with second
coordinate $b_2$, so $a_2=b_2+1$. Similarly, properness and failure of
the third guess exclude $(b_1,b_2)$ and $(b_1,b_2+1)$ for $c$; hence
$c_1=b_1+1$. Substituting these two identities into the middle guess
gives $(b_1+1+1,b_2+1+1)=b$.
The other two players may guess arbitrarily.
\end{proof}

Thus it suffices to exclude five colors. Henceforth $Q=Q_5$.

\begin{center}\small
\begin{tabular}{@{}lp{10.3cm}@{}}
\toprule
Notation&Meaning\\
\midrule
$\calC_q$&All proper colorings of the labeled cycle with palette $Q_q$.\\
$M_i^z$&The $4\times4$ indicator matrix for guessing $z$ when neither
neighbor has color $z$; row = left neighbor, column = right neighbor.\\
$s_i(z),S(z)$&The weight of this matrix and the sum of its weights over all five players.\\
$a_j$&The same five weights in the order $(1,3,5,2,4)$, in which
nonadjacent players are consecutive.\\
\bottomrule
\end{tabular}
\end{center}
The notation is defined formally below; this table fixes the two distinct
orders used in the proof and in the saved data.

\begin{lemma}[Removal of impossible guesses]\label{lem:normalization}
A winning strategy can be replaced by one satisfying
\begin{equation}\label{eq:normalization}
f_i(a,b)\notin\{a,b\}
\qquad(i\in\{1,\ldots,5\},\ a,b\in Q).
\end{equation}
Moreover, for a fixed color $z$, the replacement can preserve every
indicator $\ind[f_i(a,b)=z]$ with $a,b\ne z$.
\end{lemma}
\begin{proof}
If an answer belongs to $\{a,b\}$, it is incorrect in every proper coloring
with that observation. Replacing it by any element of
$Q\setminus\{a,b\}$ loses no correct guess. To preserve the indicated
entries, put $S=Q\setminus\{z\}$. On an observation in $S^2$, replace an
impossible answer by a color in $S\setminus\{a,b\}$; this set has at least
two elements. Neither the old nor the new answer is $z$. On observations
containing $z$, the specified indicators are not defined, so any allowed
replacement may be used. Apply this independently to all affected entries.
\end{proof}

For $z\in Q$, define
\begin{equation}\label{eq:counts}
M_i^z(a,b)=\ind[f_i(a,b)=z]\quad(a,b\in Q\setminus\{z\}),\qquad
s_i(z)=\sum_{a,b\ne z}M_i^z(a,b),\qquad
S(z)=\sum_{i=1}^5s_i(z).
\end{equation}
These counts refer to all five tables, not to the number of correct guesses
on a particular coloring. Each $s_i(z)$ is the weight of a $4\times4$
Boolean matrix and lies between 0 and~16.

\begin{lemma}[The total-count identity]\label{lem:total}
For every strategy, $\sum_zS(z)\leq125$. Under
\eqref{eq:normalization}, equality holds. Consequently, a normalized
winning strategy has a color $z$ with $S(z)\leq25$.
\end{lemma}
\begin{proof}
There are five tables, each with 25 ordered observations. An entry
$f_i(a,b)=g$ contributes to exactly one of the counts in~\eqref{eq:counts}
when $g\notin\{a,b\}$, namely to $s_i(g)$, and contributes to none
otherwise. The final assertion is the pigeonhole principle.
\end{proof}

The fixed-color preservation in Lemma~\ref{lem:normalization} has a
specific quantifier: for each chosen $z$, there is a suitable repair.
It does not assert that one repair preserves all five values $S(z)$
simultaneously. The main argument first makes any normalization and then
chooses a least-used color. The preservation clause is needed when
transferring a lower bound for a specified color back to an unnormalized
strategy, as in Corollary~\ref{cor:balance}.

The number of permitted configurations will also be useful. With $J_q$
the all-ones matrix and $I_q$ the identity matrix, proper cyclic colorings
are the closed walks of length~5 in the color-transition matrix $J_q-I_q$.
Its eigenvalues are $q-1$ once and $-1$ with multiplicity $q-1$, so
\begin{equation}\label{eq:configuration-count}
|\calC_q|=\operatorname{tr}((J_q-I_q)^5)=(q-1)^5-(q-1).
\end{equation}
In particular, $|\calC_5|=1020$ and $|\calC_4|=240$. Vertices are labeled,
so there is no division by~5. A fixed color can occur zero, one, or two
times on $C_5$; the corresponding numbers are
\begin{equation}\label{eq:layers}
240,\qquad 5\cdot4\cdot3^3=540,\qquad5\cdot4\cdot(4\cdot3)=240.
\end{equation}
Indeed, one occurrence leaves a path on four vertices, and two occurrences
must occupy one of the five nonadjacent pairs.

